\subsection{\bilingual{Proof of the Polynomial Decay of the Right Tail for \texorpdfstring{$a=0$}{a=0}}
{在 \texorpdfstring{$a=0$}{a=0} 的情形下右尾部多项式衰减的证明}}\label{section:a:0}

\bilingual{
    In this subsection, we would prove \eqref{eq:4.15} for the case of $a=0$, i.e., starting with one particle at time $0$ at position $0$. We need the following stronger moment assumption on the leftmost step $\Xi_1$ of the branching random walk $\mathcal V$.
}{
    在本小节中, 将证明 \eqref{eq:4.15} 在 $a=0$ (即从零时刻 $0$ 处的一个粒子开始) 情形下的成立. 需要对分支随机游走 $\mathcal V$ 的最左侧步长 $\Xi_1$ 做出以下更强的矩假设.
}

\setcounter{assumption}{8}
\begin{assumption}[\bilingual{Exponential Moment of Leftmost Step}{最左侧正指数矩}]
    \label{assumption:moment-exp-1}
    \bilingual{
        Under \Cref{assumption:moment-4}, we further assume that there exists some $\lambda>0$ such that
    }{
        在\Cref{assumption:moment-4} 下, 假设存在 $\lambda>0$ 使得
    }
    $\mathbb E\left[e^{\lambda\Xi_1}\right]<\infty$.
\end{assumption}

\bilingual{
    We restate the expression of \eqref{eq:4.15} for the case of $a=0$ as follows: for any $t>0$, $\zeta'>\frac{\pi^2}2\sigma^2t$, there exists some $\varepsilon=\varepsilon(t,\zeta')>0$, some $N_0=N_0(t,\zeta')>0$, such that for all $N\ge N_0$,
}{
    现将 \eqref{eq:4.15} 重述: 对任意 $t>0$, $\zeta'>\frac{\pi^2}2\sigma^2t$, 存在某个 $\varepsilon=\varepsilon(t,\zeta')>0$, 某个 $N_0=N_0(t,\zeta')>0$, 使得对于所有 $N\ge N_0$,
}
\begin{equation}
    \mathbb P\left(
        \mathcal M^N_{\ell_N}
        \ge
        \zeta'\log N
    \right)
    \le
    N^{-\varepsilon}.
    \tag{\ref{eq:4.15}}
\end{equation}

\begin{proof}
    \quad

    \bilingual{
        We are to construct a coupled killed branching random walk $\mathcal V^{\mathrm b}$ w.r.t. $\mathcal V^{N|0}$ on the same probability space such that $\mathcal Z^{\mathrm b}_n\succcurlyeq \mathcal Z^{N|0}_n$ for all $n\in[\![0,\ell_N]\!]$, that, w.h.p., $\mathcal V^{\mathrm b}$ reaches the population size of at least $N^\alpha$ for some constant $\alpha\in(0,1)$ till some random time $\tau$ of order $o_N(\log^3 N)$, and that the remaining process of $\mathcal V^{\mathrm b}$ from time $\tau$ to time $\ell_N$ is a translated $\mathcal V^{N|\alpha}$, whose minimum position at time $\ell_N$ is smaller than $\zeta'\log N$ with high probability, which has been proven above. Note that the choice of $\alpha$ should not depend on $N$, and can be arbitrary in $(0,1)$.
    }{
        下面将构造一个与 $\mathcal V^{N|0}$ 在同一概率空间上耦合的经灭杀的分支随机游走 $\mathcal V^{\mathrm b}$, 使得 $\mathcal Z^{\mathrm b}_n\succcurlyeq \mathcal Z^{N|0}_n$ 对所有 $n\in[\![0,\ell_N]\!]$ 成立, 并且 $\mathcal V^{\mathrm b}$ 在某个随机时间 $\tau$ 前以高概率达到至少 $N^\alpha$ 的粒子总数, 其中 $\alpha\in(0,1)$ 是某个不依赖于 $N$ 的常数, 并且 $\mathcal V^{\mathrm b}$ 从时刻 $\tau$ 到时刻 $\ell_N$ 的剩余过程是一个平移的 $\mathcal V^{N|\alpha}$, 其在时刻 $\ell_N$ 的最小位置以高概率小于 $\zeta'\log N$, 这已在上文证得. 注意 $\alpha$ 的选择不依赖于 $N$, 且可任取自 $(0,1)$.
    }

    \begin{figure}[htbp]
        \centering
        \begin{tikzpicture}[
            xscale=0.7,
            yscale=0.6,
            >=Stealth,
            % 样式定义
            survivor/.style={circle, fill=blue!70!black, inner sep=1.2pt},
            node-main/.style={circle, fill=blue!70!black, inner sep=1.2pt},
            node-backup/.style={circle, fill=green!60!black, inner sep=1.2pt},
            edge-main/.style={thick, blue!70!black},
            edge-backup/.style={thick, green!60!black},
            % 阴影区域样式
            fail-cloud/.style={fill=gray!20, opacity=0.6},
            success-cloud/.style={fill=blue!15, opacity=0.7},
            final-cloud/.style={fill=blue!25, opacity=0.8},
            ref-line/.style={thin, black, dashed}
        ]

            % ================= 1. 时间轴参考线 (tau_1, tau_2, T) =================
            \draw[ref-line] (1.5, -1.5) -- (1.5, 6.5) node[above] {$\tau_1$};
            \draw[ref-line] (5.0, -1.5) -- (5.0, 6.5) node[above] {$\tau_2$};
            \draw[ref-line] (10, -1.5) -- (10, 6.5) node[above] {$\tau$};
            \draw[ref-line] (17, -1.5) -- (17, 6.5) node[above] {$\ell_N$};

            % ================= 2. 演化路径: w0 到 rho_1 =================
            \node[node-main] (w0) at (0, 0) {};
            \node[left] at (w0) {$w_0$};

            \coordinate (w1) at (0.8, 0.6);
            \node[node-main] at (w1) {};
            \draw[edge-main] (w0) -- (w1);

            \node[node-main] (rho1) at (1.5, 1.2) {};
            \node[above left] at (rho1) {$\rho_1$};
            \draw[edge-main] (w1) -- (rho1);

            % ================= 3. rho_1 处的分裂与 M 标注 =================
            \node[node-main] (v1) at (2.5, 0.8) {};
            \node[node-backup] (u1) at (2.5, 2.5) {};
            \node[below] at (v1) {$v_1$};
            \node[above] at (u1) {$u_1$};

            \draw[edge-main] (rho1) -- (v1);
            \draw[edge-backup] (rho1) -- (u1);

            % M 宽度标注 (从 rho1 和 u1 所在高度引出水平线)
            \draw[thin, black!30] (1.5, 1.2) -- (3.5, 1.2); % rho1 水平线
            \draw[thin, black!30] (2.5, 2.5) -- (3.5, 2.5); % u1 水平线
            \draw[<->] (3.3, 1.2) -- (3.3, 2.5) node[midway, right] {$\le M$};

            % ================= 4. v_1 失败阴影三角 (细长, 延伸至 T 之外) =================
            \fill[fail-cloud] (2.5, 0.8) -- (6.5, 0.1) -- (6.5, 1.5) -- cycle;
            \node at (6.5, 0.8) {$\#\mathcal{Z}\notin[N^\alpha,N]$};

            % ================= 5. 重启路径: u_1 到 rho_2 =================
            \coordinate (w_mid) at (3.8, 3.2);
            \node[node-backup] at (w_mid) {};
            \draw[edge-backup] (u1) -- (w_mid);

            \node[node-backup] (rho2) at (5.0, 4.0) {};
            \node[above left] at (rho2) {$\rho_2$};
            \draw[edge-backup] (w_mid) -- (rho2);

            % ================= 6. rho_2 处的分裂与 v_2 成功三角 =================
            \node[node-main] (v2) at (6.0, 3.5) {};
            \node[node-backup] (u2) at (6.0, 5.0) {};
            \node[below] at (v2) {$v_2$};
            \node[above] at (u2) {$u_2$};

            \draw[edge-main] (rho2) -- (v2);
            \draw[edge-backup] (rho2) -- (u2);

            % v_2 成功三角 (正好延伸至 T)
            \fill[success-cloud] (6.0, 3.5) -- (10, 2.8) -- (10, 5.5) -- cycle;

            % delta log N 时间标度 (tau2 到 T)
            \draw[<->] (5.0, -1) -- (10, -1) node[midway, above] {$\ell$};

            % 最终阴影三角: 以 maxZT (10, 5.5) 为顶点
            \fill[final-cloud] (10, 5.5) -- (17, 4.0) -- (17, 7.0) -- cycle;
            \node at (14.5, 5.5) {$\mathcal V^{N|\alpha}$};

            % ================= 7. T 之后的最终粒子云 =================
            % T时刻最高点标注为顶点的三角形
            \node[survivor] (maxZT) at (10, 5.5) {};
            \node[left] at (maxZT) {$\max\mathcal Z^{\mathrm b}_\tau$};
            \node at (10, 4) {$\#\mathcal{Z}\in[N^\alpha,N]$};

            % 标注最终最小值
            \node[survivor] (minFinal) at (17, 4.0) {};
            \draw[blue!70!black, <-, thick] (17.1, 4.0) -- (17.6, 4.0) node[right] {$\mathcal{M}^{\mathrm{b}}_{\ell_N}$};

        \end{tikzpicture}
        \caption{\bilingual{Illustration of the BRW restart mechanism coupled with the \texorpdfstring{$N$-BRW}{N-BRW} when \texorpdfstring{$a=0$}{a=0}}{当 \texorpdfstring{$a=0$}{a=0} 时与 \texorpdfstring{$N$-BRW}{N-BRW} 耦合的 BRW 重启机制示意图}}
        \label{fig:coupling-restart}
    \end{figure}

    \bilingual{
        Now we specify the construction of the coupled killed branching random walk $\mathcal V^{\mathrm b}$ w.r.t. $\mathcal V^{N|0}$ (see \Cref{fig:coupling-restart}), and then we would proceed to show the aforementioned properties of $\mathcal V^{\mathrm b}$.
    }{
        现在具体说明与 $\mathcal V^{N|0}$ 耦合的经灭杀分支随机游走 $\mathcal V^{\mathrm b}$ 的构造, 参见\Cref{fig:coupling-restart}. 然后将继续证明上述 $\mathcal V^{\mathrm b}$ 的性质.
    }

    \begin{itemize}
        \item
            \bilingual{
                At time $0$ , we start with one particle at position $0$, and we let it reproduce till some time $\tau_1$ to be specified later with the offspring law
            }{
                在时刻 $0$ , 从 $0$ 处的一个粒子开始, 让它按照如下的分支律繁殖, 直到某个稍后指定的时刻 $\tau_1$为止,
            }
            \begin{equation*}
                \Xi^{(M)\mathrm b}
                :=
                \delta_{\Xi_1}
                {\bf 1}_{\{\#\Xi\ge 1\}}
                +
                \delta_{\Xi_2}
                {\bf 1}_{\{\Xi_2\le M\}}
                {\bf 1}_{\{\#\Xi\ge 2\}},
            \end{equation*}
            \bilingual{
                where $M>0$ is some constant not depending on $N$ and large enough such that $p_0:=\mathbb P(\#\Xi^{(M)\mathrm b}\ge 2)>0$. The existence of such an $M$ is guaranteed since under \Cref{assumption:psi-zero}, i.e., $\psi(0)>0$ or $\mathbb E[\#\Xi]>1$, we have $p:=\mathbb P(\#\Xi\ge 2)>0$ (or otherwise $\#\Xi\le 1$ a.s. and $\mathbb E[\#\Xi]\le 1$), and
            }{
                其中 $M>0$ 是某个不依赖于 $N$ 的常数, 足够大以使得 $p_0:=\mathbb P(\#\Xi^{(M)\mathrm b}\ge 2)>0$. 这样的 $M$ 的存在是有保证的, 因为在\Cref{assumption:psi-zero} 下, 即 $\psi(0)>0$ 或 $\mathbb E[\#\Xi]>1$, 有 $p:=\mathbb P(\#\Xi\ge 2)>0$ (否则 $\#\Xi\le 1$ a.s. 且 $\mathbb E[\#\Xi]\le 1$), 并且
            }
            \begin{equation*}
                \lim_{M\to\infty}
                \mathbb P(\#\Xi^{(M)\mathrm b}\ge 2)
                =
                \lim_{M\to\infty}
                \mathbb P(\#\Xi\ge 2,\,\Xi_2\le M)
                =
                \mathbb P(\#\Xi\ge 2)
                =
                p
                >
                0.
            \end{equation*}
            \bilingual{
                Note also that $\#\Xi^{(M)\mathrm b}\ge 1$ a.s. under \Cref{assumption:at-least-one-child}, i.e. $\#\Xi\ge 1$ a.s. And we define $\tau_1$ as the first time the population size reaches $2$ at the next generation, i.e.,
            }{
                也注意到在\Cref{assumption:at-least-one-child} 下, 即在 $\#\Xi\ge 1$ a.s. 的条件下, $\#\Xi^{(M)\mathrm b}\ge 1$ 几乎必然成立. 并定义 $\tau_1$ 为第一次使得下一代粒子总数达到 $2$ 的时刻, 即
            }
            \begin{equation*}
                \tau_1
                :=
                \inf\left\{
                    n\ge 0:\,
                    \#\mathcal Z^{\mathrm b}_{n+1}\ge 2
                \right\}.
            \end{equation*}
            \bilingual{
                This $\tau_1$ is generally not a stopping time for the generation filtration. Indeed, $\mathcal F_0=\{\emptyset,\Omega\}$, whereas $\{\tau_1\le0\}=\{\#\Xi^{(M)\mathrm b}_{\emptyset}\ge2\}$ has probability $p_0$. In the allowed case $0<p_0<1$, this event is not in $\mathcal F_0$. The first observable split generation is $\sigma_1:=\tau_1+1$.
            }{
                对逐代域流而言, 这个 $\tau_1$ 一般不是停时. 事实上 $\mathcal F_0=\{\emptyset,\Omega\}$, 但 $\{\tau_1\le0\}=\{\#\Xi^{(M)\mathrm b}_{\emptyset}\ge2\}$ 的概率为 $p_0$. 在允许出现的 $0<p_0<1$ 情形下, 此事件不属于 $\mathcal F_0$. 首次可观测的分叉代为 $\sigma_1:=\tau_1+1$.
            }
            \bilingual{
                Note that $\tau_1$ does not depend on $N$, and $\tau_1<\infty$ a.s. since $p_0>0$. We labeled the particles from time $0$ to time $\tau_1$ as $w_n\in\mathbb V$, $n\in[\![0,\tau_1]\!]$, where $w_0$ is the ancestor at time $0$, and for each $n\in[\![0,\tau_1-1]\!]$, $w_{n+1}$ is the unique child of $w_n$. We also denote $\rho_1:=w_{\tau_1}$, $v_1$ as the left-most child of $\rho_1$, and $u_1$ as the second to the left-most child of $\rho_1$.
            }{
                注意 $\tau_1$ 不依赖于 $N$, 且 $\tau_1<\infty$ 几乎必然成立, 此因 $p_0>0$. 将从时刻 $0$ 到时刻 $\tau_1$ 的粒子记为 $w_n\in\mathbb V$, $n\in[\![0,\tau_1]\!]$, 其中 $w_0$ 是零时刻的祖先, 对于每个 $n\in[\![0,\tau_1-1]\!]$, $w_{n+1}$ 是 $w_n$ 的唯一子代. 还记 $\rho_1:=w_{\tau_1}$, $v_1$ 是 $\rho_1$ 的最左侧的子代, 而 $u_1$ 是 $\rho_1$ 的第二最左侧的子代.
            }
        \item
            \bilingual{
                From time $\tau_1+1$ onward, suppressing the reproduction of $u_1$, we start with one particle at position $\mathcal V^{\mathrm b}(v_1)$, and we let it reproduce according to the offspring law $\Xi^{(M)}$ till some time $\tau_1+\ell$, where $\ell:=\lfloor\frac 1{\log m}\alpha\log N\rfloor\ge 2$, and
            }{
                从时刻 $\tau_1+1$ 开始, 抑制 $u_1$ 的繁殖, 并从位置 $\mathcal V^{\mathrm b}(v_1)$ 的一个粒子开始, 让它按照如下的分支律 $\Xi^{(M)}$ 繁殖, 直到某时刻 $\tau_1+\ell$, 其中 $\ell:=\lfloor\frac 1{\log m}\alpha\log N\rfloor\ge 2$, 并且
            }
            \begin{equation*}
                \Xi^{(M)}
                :=
                \delta_{\Xi_1}
                {\bf 1}_{\{\Xi_1\le M\}}
                {\bf 1}_{\{\#\Xi\ge 1\}}
                +
                \delta_{\Xi_2}
                {\bf 1}_{\{\Xi_2\le M\}}
                {\bf 1}_{\{\#\Xi\ge 2\}},
            \end{equation*}
            \bilingual{
                where $M$ is the same constant as above, but with an even stronger condition that $m:=\mathbb E[\#\Xi^{(M)}]>1$. We will fix the choice of $M$ from now on. The existence of such an $M$ is guaranteed since under \Cref{assumption:psi-zero}, i.e., $\psi(0)>0$ or $\mathbb E[\#\Xi]>1$, we have $\mathbb P(\#\Xi\ge 2)=p>0$ and we invoke the reasoning as above and moreover, under \Cref{assumption:at-least-one-child}, i.e. $\#\Xi\ge 1$ a.s., we have
            }{
                其中 $M$ 是与前述相同的常数, 但需要更强的性质 $m:=\mathbb E[\#\Xi^{(M)}]>1$. 从现在起固定 $M$ 的选择. 该 $M$ 的存在性是有保证的, 此因在\Cref{assumption:psi-zero} 下, 即 $\psi(0)>0$ 或 $\mathbb E[\#\Xi]>1$, 有 $\mathbb P(\#\Xi\ge 2)=p>0$ 并且可以使用前述相同的推理, 而且在\Cref{assumption:at-least-one-child} 下, 即 $\#\Xi\ge 1$ a.s., 有
            }
            \begin{align*}
                \lim_{M\to\infty}
                \mathbb E\left[\#\Xi^{(M)}\right]
                &
                =
                \lim_{M\to\infty}
                \mathbb E\left[
                    {\bf 1}_{\{\Xi_1\le M\}}
                    {\bf 1}_{\{\#\Xi\ge 1\}}
                    +
                    {\bf 1}_{\{\Xi_2\le M\}}
                    {\bf 1}_{\{\#\Xi\ge 2\}}
                \right]
                \\ &
                =
                \mathbb P(\#\Xi\ge 1)
                +
                \mathbb P(\#\Xi\ge 2)
                =
                1+p
                >
                1.
            \end{align*}
            \bilingual{
                Note that we only have $\mathbb E[\#\Xi^{(M)\mathrm b}]>1$ but not $\#\Xi^{(M)\mathrm b}\ge 1$ a.s., which implies the process may get extinct with positive probability from time $\tau_1+1$ till $\tau_1+\ell$. If this happens, we start over from time $\tau_1+1$ at position $u_1$, i.e., the second to the left-most child of $\rho_1$, and repeat the same procedure as above, see the items below for details. Note also that since $\alpha$, $M$, and $m$ does not depend on $N$, we can always let $N$ large enough such that $\ell\ge 2$.
            }{
                注意只有 $\mathbb E[\#\Xi^{(M)\mathrm b}]>1$ 而没有 $\#\Xi^{(M)\mathrm b}\ge 1$ a.s., 这意味着从时刻 $\tau_1+1$ 到 $\tau_1+\ell$ 的过程中可能以正概率灭绝. 若发生此种情况, 就从时刻 $\tau_1+1$ 的位置 $u_1$ (即 $\rho_1$ 的第二最左侧的子代) 重新开始, 并重复上述相同的过程, 详情见下面的条目. 还注意到既然 $\alpha$, $M$, 和 $m$ 不依赖于 $N$, 可以总是让 $N$ 足够大以使得 $\ell\ge 2$.
            }
        \item
            \bilingual{
                At time $\tau_1+\ell$, if we reach the population size of at least $N^\alpha$ (without extinction) and not exceeding $N$, we start with $\lfloor N^\alpha\rfloor$ particles at position $\max\mathcal Z^{\mathrm b}_\tau$ where we set $\tau:=\tau_1+\ell$, and we let them reproduce according to the original offspring law $\Xi$ and the selection mechanism of keeping the $N$ left-most particles till time $\ell_N$.
            }{
                在时刻 $\tau_1+\ell$, 如果达到至少 $N^\alpha$ 的粒子总数 (没有灭绝) 且不超过 $N$, 就从位置 $\max\mathcal Z^{\mathrm b}_\tau$ 的 $\lfloor N^\alpha\rfloor$ 个粒子开始, 其中 $\tau:=\tau_1+\ell$, 让它们按照原始的分支律 $\Xi$ 和保留最左侧 $N$ 个粒子的选择机制繁殖直到时刻 $\ell_N$.
            }
        \item
            \bilingual{
                Otherwise, if we do not reach the population size of $N^\alpha$ at time $\tau_1+\ell$, or if we get extinct or get to the population of $N$ at some time between $\tau_1+1$ and $\tau_1+\ell$, we start over from time $\tau_1+1$ at position $u_1$, i.e., the second to the left-most child of $\rho_1$, and repeat the same procedure as above, i.e., we let it reproduce according to the offspring law $\Xi^{(M)b}$ till time $\tau_2$, the first time the population size of the descendants of $u_1$ reaches $2$ at the next generation,
            }{
                否则, 如果在时刻 $\tau_1+\ell$ 没有达到粒子总数 $N^\alpha$, 或者, 如果在 $\tau_1+1$ 和 $\tau_1+\ell$ 之间的某个时刻灭绝或达到粒子总数$N$, 就从时刻 $\tau_1+1$ 的位置 $u_1$ (即 $\rho_1$ 的第二最左侧的子代) 重新开始, 并重复上述相同的过程, 即让它按照分支律 $\Xi^{(M)b}$ 繁殖直到时刻 $\tau_2$, 其中 $\tau_2$ 是第一次使得 $u_1$ 的后代在下一代达到 $2$ 的时刻,
            }
            \bilingual{
                In the defining display below, $\widehat{\mathcal Z}^{\mathrm b(u_1)}$ denotes the hypothetical reserve-subtree process generated from the pre-sampled offspring marks of $u_1$ on every outcome, including when trial 1 succeeds.
            }{
                在下述定义式中, $\widehat{\mathcal Z}^{\mathrm b(u_1)}$ 表示由 $u_1$ 的预采样子代标记在所有样本结果上构造的候选备用子树过程, 包括第一次试验成功的情形.
            }
            \begin{equation*}
                \tau_2
                :=
                \inf\left\{
                    n\ge \tau_1+1:\,
                    \#\widehat{\mathcal Z}^{\mathrm b(u_1)}_{n+1}\ge 2
                \right\},
            \end{equation*}
            \bilingual{
                For the stopping-time argument, this is an \emph{auxiliary candidate time} defined from the pre-sampled offspring marks of the reserve subtree rooted at $u_1$, whether or not trial 1 succeeds. Only the actual use of this reserve subtree depends on trial 1 failing. We make the same distinction for every later candidate.
            }{
                为论证停时性质, 这里的 $\tau_2$ 是根据预采样的备用子树 $u_1$ 内的分支标记定义的\emph{辅助候选时刻}, 无论第一次试验是否成功都定义它. 只有是否实际使用该备用子树才取决于第一次试验失败. 对后续每个候选也作相同区分.
            }
            \bilingual{
                where we use the superscript $u_1$ to denote the descendants of $u_1$, with the labeling $w_n\in \mathbb V$, $n\in[\![\tau_1+1,\tau_2]\!]$, $\rho_2$, $v_2$ and $u_2$ alike, and then we let $v_2$ reproduce according to the offspring law $\Xi^{(M)}$ for $(\ell-1)$ generations till time $\tau_2+\ell>\tau_1+\ell$. Note that during this procedure, we suppress the reproduction of $v_1$, i.e., the left-most child of $\rho_1$. Note also that $\tau_2-\tau_1$ does not depend on $N$, and $\tau_2<\infty$ a.s. since $p_0>0$.
            }{
                其中使用上标 $u_1$ 来表示 $u_1$ 的后代, 粒子标记 $w_n\in \mathbb V$, $n\in[\![\tau_1+1,\tau_2]\!]$, $\rho_2$, $v_2$ 和 $u_2$ 类似, 然后让 $v_2$ 按照分支律 $\Xi^{(M)}$ 繁殖 $(\ell-1)$ 代直到时刻 $\tau_2+\ell>\tau_1+\ell$. 注意在此过程中, 需要抑制 $\rho_1$ 的最左侧的子代 $v_1$ 的繁殖. 还需注意 $\tau_2-\tau_1$ 不依赖于 $N$, 且 $\tau_2<\infty$ 几乎必然成立, 此因 $p_0>0$.
            }
        \item
            \bilingual{
                On the pre-sampled marked tree, define every reserve lineage and its candidate split time $\sigma_k:=\tau_k+1$ recursively, irrespective of the outcomes of preceding growth trials. At $T_k:=\sigma_k+\ell-1=\tau_k+\ell$, let $A_k$ be the candidate event that the descendants of $v_k$ reach a population in $[N^\alpha,N]$ without an earlier forbidden population size. Set $\kappa:=\inf\{k\ge1:A_k\}$ and $\tau:=T_\kappa$, with $\tau=\infty$ if no trial succeeds. Actual use of the $k$-th candidate is restricted to $A_1^{\mathrm c}\cap\cdots\cap A_{k-1}^{\mathrm c}$. To apply the branching property at $\tau$, one must prove that each $\sigma_k$ is an $(\mathcal F_n)$-stopping time and that $A_k$ is observable at $T_k$. Then $\{\tau\le n\}$ is the union of $\{T_k\le n,A_k,A_1^{\mathrm c},\ldots,A_{k-1}^{\mathrm c}\}$ over $k\ge1$, hence $\tau$ is a stopping time. These model-specific measurability facts are still to be verified. We will show that, for some $\delta>0$ and $\gamma>0$ to be specified later, $\kappa\le K:=\lfloor\delta\log N\rfloor$ and $\tau\le T:=\gamma\log N$ holds with high probability. Then, at time $\tau$, on the event $\{\tau<\infty\}$, we start with $\lfloor N^\alpha\rfloor$ particles at position $\max\mathcal Z^{\mathrm b}_\tau$, and let them reproduce according to the original offspring law $\Xi$ and the selection mechanism of keeping the $N$ left-most particles till time $\ell_N$.
            }{
                在预采样标记树上, 不依赖前次试验成败地递归定义所有备用谱系及其候选分叉代 $\sigma_k:=\tau_k+1$. 在 $T_k:=\sigma_k+\ell-1=\tau_k+\ell$ 定义候选成功事件 $A_k$: $v_k$ 的后代在未提前达到禁止的粒子数的前提下, 达到区间 $[N^\alpha,N]$ 内的粒子数. 令 $\kappa:=\inf\{k\ge1:A_k\}$ 及 $\tau:=T_\kappa$; 若没有成功的试验则令 $\tau=\infty$. 第 $k$ 个候选仅在 $A_1^{\mathrm c}\cap\cdots\cap A_{k-1}^{\mathrm c}$ 上实际启用. 要在 $\tau$ 处应用分支性质, 需证明每个 $\sigma_k$ 是 $(\mathcal F_n)$-停时, 且 $A_k$ 在 $T_k$ 可观测. 此时 $\{\tau\le n\}$ 是 $\{T_k\le n,A_k,A_1^{\mathrm c},\ldots,A_{k-1}^{\mathrm c}\}$ 对 $k\ge1$ 的并, 因而 $\tau$ 是停时. 这些模型特定的可测性仍待验证. 下面将证明对于某个稍后指定的 $\delta>0$ 和 $\gamma>0$, 以高概率成立 $$\kappa\le K:=\lfloor\delta\log N\rfloor,\quad \tau\le T:=\gamma\log N.$$ 然后在时刻 $\tau$, 在事件 $\{\tau<\infty\}$ 上, 从位置 $\max\mathcal Z^{\mathrm b}_\tau$ 的 $\lfloor N^\alpha\rfloor$ 个粒子开始, 让它们按照原始的分支律 $\Xi$ 和保留最左侧 $N$ 个粒子的选择机制繁殖直到时刻 $\ell_N$.
            }
    \end{itemize}
    \bilingual{
        The above exploration is at most bi-branching before $\tau$ and keeps no more than $N$ particles. However, after a failed trial it retrospectively resumes the sibling $u_k$ from an earlier generation. Consequently the displayed generation process $\mathcal Z^{\mathrm b}_n$ need not be adapted to the generation filtration $\mathcal F_n$, as required by \Cref{lemma:4.8}. A separate causal coupling or another comparison proof is needed to justify $\mathcal Z^{\mathrm b}_n\succcurlyeq\mathcal Z^{N|0}_n$. Simply waiting until each failed trial ends before starting the next one would use $K\ell=O((\log N)^2)$ generations when $K$ and $\ell$ are both of logarithmic order; it does not preserve the present $O(\log N)$ restart-time estimate. One possible repair is to retain a causal backbone and run reserve trials concurrently from siblings born along it, with a population cap and a separate spatial estimate. That replacement has not been proved here. The following bound is conditional on the missing comparison:
    }{
        上述探索在 $\tau$ 前至多二分支, 保留的粒子数也不超过 $N$. 但是一次试验失败后, 构造会从较早一代的兄弟 $u_k$ 回溯继续, 因而所写的逐代过程 $\mathcal Z^{\mathrm b}_n$ 未必适应逐代域流 $\mathcal F_n$, 而\Cref{lemma:4.8} 明确要求这种适应性. 要证明 $\mathcal Z^{\mathrm b}_n\succcurlyeq\mathcal Z^{N|0}_n$, 需要另建因果耦合或给出其他比较证明. 若只在每次失败后从当代重新开始, 当 $K$ 与 $\ell$ 均为对数阶时, 共需 $K\ell=O((\log N)^2)$ 代, 不再保有当前 $O(\log N)$ 的重启时间估计. 一种可能的修正是保留逐代演化的主脊, 让沿途兄弟支系同时尝试增长, 并另外控制总人口与空间位移. 本文尚未证明这一替代构造. 以下界以缺失的比较为前提:
    }
    \begin{align}
        &\quad
        \mathbb P\left(
            \mathcal M^{N|0}_{\ell_N}
            \ge
            \zeta'\log N
        \right)
        % \\
        % &
        \le
        \mathbb P\left(
            \mathcal M^{\mathrm b}_{\ell_N}
            \ge
            \zeta'\log N
        \right)
        \nonumber
        \\
        &
        \le
        \mathbb P\left(
            \kappa>K
        \right)
        +
        \mathbb P\left(
            \kappa\le K,\,
            \tau>T
        \right)
        \nonumber \\ &\quad
        +
        \mathbb P\left(
            \mathcal M^{\mathrm b}_{\ell_N}
            \ge
            \zeta'\log N,\,\kappa\le K,\, \tau\le T
        \right).
        \label{eq:a:0:1}
    \end{align}

    \bilingual{
        For the first term in the last equation, the observable first-bifurcation generation is $\sigma_1:=\tau_1+1$, rather than $\tau_1$: the latter is defined by looking at the next generation and need not be a stopping time. More importantly, conditioning on the full generation $\sigma$-algebra $\mathcal F_{\tau_1+\ell}$ would reveal part of the unused sibling subtree rooted at $u_1$ and would not make the next trial fresh. Let $\mathscr H_j$ denote the information revealed by the exploration through the end of trial $j$, excluding reproduction marks in the still-unused sibling subtrees. The factorization below requires proving that each measurably chosen unused reserve subtree is independent of $\mathscr H_j$. One may prove this directly by partitioning over its countably many possible labelled roots and using independence of disjoint pre-sampled marks; a general stopping-line theorem is sufficient but not necessary. The selected root must be determined without inspecting marks inside that reserve subtree. These exploration measurability and independence claims have not yet been verified for the present construction.
    }{
        对于上式中的第一项, 可观测的首次分叉代是 $\sigma_1:=\tau_1+1$, 而不是 $\tau_1$: 后者按下一代的粒子数定义, 一般未必是停时. 更重要的是, 若条件化于整代信息 $\mathcal F_{\tau_1+\ell}$, 就已揭示备用兄弟 $u_1$ 子树中较早代的部分信息, 下一次试验不再相对它是全新的. 记 $\mathscr H_j$ 为探索到第 $j$ 次试验结束时已揭示的信息, 不包含尚未使用的兄弟子树内的繁殖标记. 下面的乘积分解须证明可测选出的每棵未使用备用子树独立于 $\mathscr H_j$. 可对其可数个可能的带标记根分割, 再利用预采样树中互不相交坐标的独立性直接证明; 通用停止线定理是充分工具, 但并非必需. 选根时不能查看该备用子树内部的标记. 本构造的探索可测性和独立性尚未验证.
    }
    \begin{align*}
        \mathbb P\left(
            \kappa>K
        \right)
        &
        =
        \mathbb P\left(
            \kappa>K,\,
            \forall k\in[\![1,K]\!]:\,
            \#\mathcal Z^{\mathrm b}_{\tau_k+\ell}
            \notin[N^\alpha,N]
        \right)
        \\ &
        =
        \mathbb E\left[
            \mathbb E\left(
                {\bf 1}_{\{\kappa>K\}}
                \prod_{k=1}^K
                {\bf 1}{
                    \left\{
                        \#\mathcal Z^{\mathrm b}_{\tau_k+\ell}
                        \notin[N^\alpha,N]
                    \right\}
                }
                \middle|
                \mathscr H_1
            \right)
        \right]
        \\ &
        =
        \mathbb E\left[
            {\bf 1}_{\{
                \#\mathcal Z^{\mathrm b}_{\tau_1+\ell}
                \notin[N^\alpha,N]
            \}}
            \mathbb E\left(
                {\bf 1}_{\{\kappa>K\}}
                \prod_{k=2}^K
                {\bf 1}{
                    \left\{
                        \#\mathcal Z^{\mathrm b}_{\tau_k+\ell}
                        \notin[N^\alpha,N]
                    \right\}
                }
                \middle|
                \mathscr H_1
            \right)
        \right]
        \\ &
        =
        \left(
            \mathbb P\left(
                \#\mathcal Z^{\mathrm b}_{\tau_1+\ell}
                \notin[N^\alpha,N]
            \right)
        \right)^K,
    \end{align*}
    \bilingual{
        where
    }{
        其中
    }
    \begin{align*}
        \mathbb P\left(
            \#\mathcal Z^{\mathrm b}_{\tau_1+\ell}
            \notin[N^\alpha,N]
        \right)
        &
        =
        \mathbb E\left[
            \mathbb E\left(
                {\bf 1}_{
                    \left\{
                        \#\mathcal Z^{\mathrm b}_{\tau_1+\ell}
                        \notin[N^\alpha,N]
                    \right\}
                }
                \middle|
                \mathscr F_{\tau_1+1}
            \right)
        \right]
        \\ &
        =
        \mathbb E\left[
            {\bf 1}_{
                \left\{
                    \#\mathcal Z^{(M)}_{\ell-1}
                    \notin[N^\alpha,N]
                \right\}
            }
        \right]
        % \\ &
        =
        1-\mathbb P\left(
            \#\mathcal Z^{(M)}_{\ell-1}
            \in[N^\alpha,N]
        \right),
    \end{align*}
    \bilingual{
        and $\mathcal Z^{(M)}$ is the killed branching random walk with offspring law $\Xi^{(M)}$ and starting with one particle at position $0$. As a standard result for supercritical Galton-Watson process, we have, for $m:=\mathbb E[\#\Xi^{(M)}]>1$, the martingale $({\#\mathcal Z^{(M)}_n}/{m^n})_{n\ge 0}$ with expection $1$ converges in distribution to some random variable $W$ which, by Kesten--Stigum Theorem, is not degenerate, i.e., $\mathbb E[W]=1$, under the ``LlogL'' condition $\mathbb E[\#\Xi^{(M)}\log \#\Xi^{(M)}]<\infty$, which is guaranteed by definition of $\Xi^{(M)}$ that $\#\Xi^{(M)}\le 2$ being at most bi-branching. Note that if we were to use only the truncation to define $\Xi^{(M)}$ and let it reproduce in infinite numbers, we might not overcome this necessary condition. Hence, we have
    }{
        并且其中的 $\mathcal Z^{(M)}$ 是如下的分支随机游走, 其分支律为 $\Xi^{(M)}$ 并且从位置 $0$ 的一个粒子开始. 作为超临界 Galton-Watson 过程的一个标准结果, 对于 $m:=\mathbb E[\#\Xi^{(M)}]>1$, 均值为$1$的鞅 $({\#\mathcal Z^{(M)}_n}/{m^n})_{n\ge 0}$ 依分布收敛于某个随机变量 $W$, 根据 Kesten--Stigum 定理, 在 ``LlogL'' 条件 $\mathbb E[\#\Xi^{(M)}\log \#\Xi^{(M)}]<\infty$ 下, $W$ 非退化, 即 $\mathbb E[W]=1$, 该条件由 $\Xi^{(M)}$ 的定义所保证, 此因 $\#\Xi^{(M)}\le 2$ 是至多二分支的. 注意如果我们仅使用截断来定义 $\Xi^{(M)}$ 并让它无限繁殖, 就可能无法达成这个必要条件. 因此, 有
    }
    \begin{align*}
        &\quad
        \liminf_{N\to\infty}
        \mathbb P\left(
            \#\mathcal Z^{(M)}_{\ell-1}
            \in[N^\alpha,N]
        \right)
        \\
        &
        =
        \liminf_{N\to\infty}
        \left(
            \mathbb P\left(
                \#\mathcal Z^{(M)}_{\ell-1}
                \ge
                N^\alpha
            \right)
            -
            \mathbb P\left(
                \#\mathcal Z^{(M)}_{\ell-1}
                >N
            \right)
        \right)
        \\ &
        \ge
        \liminf_{N\to\infty}
        \mathbb P\left(
            \#\mathcal Z^{(M)}_{\ell-1}
            \ge
            N^\alpha
        \right)
        -
        \liminf_{N\to\infty}
        \mathbb P\left(
            \#\mathcal Z^{(M)}_{\ell-1}
            >N
        \right)
    \end{align*}
    \bilingual{
        where, by Markov's inequality,
    }{
        其中, 根据 Markov 不等式,
    }
    \begin{align*}
        \mathbb P\left(
            \#\mathcal Z^{(M)}_{\ell-1}
            >N
        \right)
        &
        \le
        \mathbb P\left(
            \frac{\#\mathcal Z^{(M)}_{\ell-1}}{m^{\ell-1}}
            \ge
            N^{1-\alpha}
        \right)
        % \\ &
        % \le
        % \frac{\mathbb E\left[
        %     \#\mathcal Z^{(M)}_{\ell-1}
        % \right]}{m^{\ell-1}N^{1-\alpha}}
        =
        \frac{m^{\ell-1}}{m^{\ell-1}N^{1-\alpha}}
        =
        N^{\alpha-1}
        \to
        0,
    \end{align*}
    \bilingual{
        when $N\to\infty$, and thus
    }{
        当 $N\to\infty$ 时成立, 从而
    }
    \begin{align*}
        \liminf_{N\to\infty}
        \mathbb P\left(
            \#\mathcal Z^{(M)}_{\ell-1}
            \in[N^\alpha,N]
        \right)
        &
        \ge
        \liminf_{N\to\infty}
        \mathbb P\left(
            \frac{\#\mathcal Z^{(M)}_{\ell-1}}{m^{\ell-1}}
            \ge
            1
        \right)
        -
        0
        \\ &
        \ge
        \mathbb P(W\ge 1)
        =
        2p_1
        >
        0,
    \end{align*}
    \bilingual{
        where the last inequality holds since otherwise we have $W<1$ a.s. and $\mathbb E[W]<1$, which contradicts with $\mathbb E[W]=1$ obtained by Kesten--Stigum Theorem for $m>1$. Hence, for some $N_1=N_1(\alpha,\delta)>0$ and all $N\ge N_1$, we have $\ell\ge 2$ and
    }{
        其中最后一个不等式成立, 否则就有 $W<1$ a.s. 且 $\mathbb E[W]<1$, 这与 Kesten--Stigum 定理对于 $m>1$ 得到的 $\mathbb E[W]=1$ 矛盾. 因此, 对于某个 $N_1=N_1(\alpha,\delta)>0$ 和所有 $N\ge N_1$, 有 $\ell\ge 2$ 且
    }
    % \begin{equation*}
    $
        \mathbb P\left(
            \#\mathcal Z^{(M)}_{\ell-1}
            \in[N^\alpha,N]
        \right)
        \ge
        p_1,
    % \end{equation*}
    $
    \bilingual{
        and therefore, for $\varepsilon_1=p_1\delta>0$ and all $N\ge N_1$,
    }{
        因此, 对于$\varepsilon_1=p_1\delta>0$ 和所有 $N\ge N_1$,
    }
    \begin{align}
        \mathbb P\left(
            \kappa>K
        \right)
        &
        =
        \left(
            1
            -
            \mathbb P\left(
                \#\mathcal Z^{\mathrm b}_{\tau_1+\ell}
                \in[N^\alpha,N]
            \right)
        \right)^K
        \le
        (1-p_1)^{\delta\log N}
        \nonumber
        \\
        &
        \le
        (e^{-p_1})^{\delta\log N}
        =
        N^{-p_1\delta}
        =:
        N^{-\varepsilon_1}.
        \label{eq:kappa:bound}
    \end{align}

    \bilingual{
        For the second term of \eqref{eq:a:0:1}, since $\tau=\tau_\kappa+\ell$,
    }{
        对于 \eqref{eq:a:0:1} 中的第二项, 既然 $\tau=\tau_\kappa+\ell$,
    }
    \begin{align*}
        % &\quad
        \mathbb P\left(
            \kappa\le K,\, \tau>T
        \right)
        % &
        \le
        \mathbb P\left(
            \tau_K+\ell>T
        \right)
        % \\
        % &
        =
        \mathbb P\left(
            \sum_{j=1}^K
            G_j
            +
            (K-1)
            +
            \ell
            >
            T
        \right),
    \end{align*}
    \bilingual{
        where $G_j=\tau_j-\tau_{j-1}-1$, $j\in\mathbb N$ (where $\tau_0:=-1$ and we extend the definition for $j>\kappa$), are i.i.d. geometric random variables with success probability $p_0=\mathbb P(\#\Xi^{(M)\mathrm b}\ge 2)>0$ and mean $\frac 1{p_0}-1$. Note that there exists some $0<\theta<-\log(1-p_0)$ such that $\mathbb E[e^{\theta G_1}]<\infty$. Indeed,
    }{
        其中 $G_j=\tau_j-\tau_{j-1}-1$, $j\in\mathbb N$ (其中 $\tau_0:=-1$ 且对于 $j>\kappa$ 也扩展定义), 是成功概率为 $p_0=\mathbb P(\#\Xi^{(M)\mathrm b}\ge 2)>0$ 的独立同分布几何随机变量列, 均值为 $\frac 1{p_0}-1$. 注意到存在某个 $0<\theta<-\log(1-p_0)$ 使得 $\mathbb E[e^{\theta G_1}]<\infty$. 事实上,
    }
    \begin{align*}
        \mathbb E[e^{\theta G_1}]
        &
        =
        \sum_{k=0}^\infty
        e^{\theta k}
        (1-p_0)^k
        p_0
        =
        \frac{p_0}{1-(1-p_0)e^\theta}
        <
        \infty.
    \end{align*}
    \bilingual{
        By Markov's inequality, we have
    }{
        根据 Markov 不等式, 有
    }
    \begin{align}
        \mathbb P\left(
            \kappa\le K,\, \tau>T
        \right)
        &
        \le
        \mathbb P\left(
            e^{\theta\sum_{j=1}^K G_j}
            \ge
            e^{\theta(T-K-\ell+1)}
        \right)
        \le
        \frac{\mathbb E\left[
            e^{\theta G_1}
        \right]^K}{e^{\theta(T-K-\ell)}}
        \nonumber \\ &
        \le
        N^{
            \theta\left(
                \delta+\frac\alpha{\log m}
                -
                \gamma
            \right)
            +
            \delta\log\mathbb E[e^{\theta G_1}]
        }
        =:
        N^{-\varepsilon_2},
        \label{eq:kappa:tau:bound}
    \end{align}
    \bilingual{
        where, given $\Xi$ (and thus $p_0$, $m$, law of $G_1$ and $\theta$ are all determined), for $\delta>0$ and $\alpha>0$ to be specified later, we fix $\gamma>0$ such that $\varepsilon_2>0$, i.e.,
    }{
        其中, 给定 $\Xi$ (因此 $p_0$, $m$, $G_1$的分布律以及 $\theta$ 都被确定), 对于稍后指定的 $\delta>0$ 和 $\alpha>0$, 固定 $\gamma>0$ 使得 $\varepsilon_2>0$, 即
    }
    \begin{equation}
        \gamma
        >
        \frac\delta\theta
        \mathbb E\left[
            e^{\theta G_1}
        \right]
        +
        \delta
        +
        \frac\alpha{\log m}.
        \label{eq:gamma:condition}
    \end{equation}

    \bilingual{
        For the third term of \eqref{eq:a:0:1}, by the triangular inequality, for some $c>0$ to be specified later, and
    }{
        对于 \eqref{eq:a:0:1} 中的第三项, 根据三角不等式, 对于某个稍后指定的 $c>0$, 和
    }
    % \begin{equation*}
    $
        x_N
        :=
        c\delta\log N
        +
        (\kappa-1)M
        +
        \left(
            \max\mathcal Z^{\mathrm b}_\tau
            - \mathcal V^{\mathrm b}(\rho_\kappa)
        \right),
    % \end{equation*}
    $
    \bilingual{
        we have
    }{
        成立
    }
    \begin{align}
        &\quad
        \mathbb P\left(
            \mathcal M^{\mathrm b}_{\ell_N}
            \ge
            \zeta'\log N,\,\kappa\le K,\, \tau\le T
        \right)
        \nonumber
        \\
        &
        =
        \mathbb P\left(
            \max\mathcal Z^{\mathrm b}_\tau
            +
            \left(
                \mathcal M^{\mathrm b}_{\ell_N}
                -
                \max\mathcal Z^{\mathrm b}_\tau
            \right)
            \ge
            x_N+(\zeta'\log N-x_N),\,
            \kappa\le K,\, \tau\le T
        \right)
        \nonumber
        \\ &
        \le
        \mathbb P\left(
            \max\mathcal Z^{\mathrm b}_\tau
            \ge
            x_N,\,
            \kappa\le K,\,\tau\le T
        \right)
        \nonumber \\ &\quad
        +
        \mathbb P\left(
            \mathcal M^{\mathrm b}_{\ell_N}
            -
            \max\mathcal Z^{\mathrm b}_\tau
            \ge
            \zeta'\log N-x_N,\,
            \kappa\le K,\,\tau\le T
        \right),
        \label{eq:a:0:2}
    \end{align}
    \bilingual{
        For the first part of the last equation, we denote $\tau_0:=-1$ as above and
    }{
        对最后一个等式的第一部分, 记 $\tau_0:=-1$ 如上, 且
    }
    \begin{equation*}
        X_k
        :=
        \sum_{i=\tau_{k-1}+1}^{\tau_k-1}
        \left(
            \mathcal V^{\mathrm b}(w_{i+1})
            -
            \mathcal V^{\mathrm b}(w_i)
        \right),
        \quad
        S_k
        :=
        \sum_{i=1}^k
        X_i,
        \quad
        k\in[\![1,\kappa]\!].
    \end{equation*}
    \bilingual{
        and since
    }{
        且既然
    }
    \begin{align}
        \max\mathcal Z^{\mathrm b}_\tau
        &
        =
        S_\kappa
        +
        \sum_{k=1}^{\kappa-1}
        \left(
            \mathcal V^{\mathrm b}(u_k)
            -
            \mathcal V^{\mathrm b}(\rho_k)
        \right)
        +
        \left(
            \max\mathcal Z^{\mathrm b}_\tau
            -
            \mathcal V^{\mathrm b}(\rho_\kappa)
        \right)
        \nonumber \\ &
        \le
        S_\kappa
        +
        (\kappa-1)M
        +
        \left(
            \max\mathcal Z^{\mathrm b}_\tau
            -
            \mathcal V^{\mathrm b}(\rho_\kappa)
        \right),
        \label{eq:a:0:max-Zb-tau}
    \end{align}
    \bilingual{
        we have
    }{
        因而有
    }
    \begin{align*}
        &\quad
        \mathbb P\!\left(
            \max\mathcal Z^{\mathrm b}_\tau
            \ge
            x_N,\,\kappa\le K,\,\tau\le T
        \right)
        \\
        &
        \le
        \mathbb P\left(
            S_\kappa
            \ge
            c\delta\log N,\,\kappa\le K,\,\tau\le T
        \right)
        \\ &
        \le
        \mathbb P\!\left(
            \!
            \max_{1\le k\le K}
            \!
            S_k
            \ge
            c\delta\log N
            \!\!
        \right)
        % \\ &
        =
        \mathbb P\left(
            \!
            \max_{1\le k\le K}
            \!
            e^{\lambda S_k}
            \ge
            e^{\lambda c\delta\log N}
            \!
        \right)
    \end{align*}
    \bilingual{
        Let $B:=\{\#\Xi^{(M)\mathrm b}\ge2\}$ be the split event for one reproduction, so $p_0=\mathbb P(B)$. The waiting time $G_k$ counts consecutive failures of $B$. A displacement $\Xi_1$ observed during this wait is coupled to $B^{\mathrm c}$ and does not have the unconditional law of $\Xi_1$. Define, for a sufficiently small $\lambda>0$ allowed by \Cref{assumption:moment-exp-1},
    }{
        令 $B:=\{\#\Xi^{(M)\mathrm b}\ge2\}$ 为一次繁殖中的分叉事件, 则 $p_0=\mathbb P(B)$. 等待时间 $G_k$ 计数连续的 $B$ 失败次数. 在等待期间观测到的位移 $\Xi_1$ 与 $B^{\mathrm c}$ 耦合, 不服从无条件的 $\Xi_1$ 分布. 在\Cref{assumption:moment-exp-1} 允许的范围内取足够小的 $\lambda>0$, 定义
    }
    \begin{equation*}
        A_\lambda:=\mathbb E\left[e^{\lambda\Xi_1}{\bf1}_{B^{\mathrm c}}\right].
    \end{equation*}
    \bilingual{
        We have $A_0=1-p_0<1$. Dominated convergence, using $\mathbb E[e^{\lambda_0\Xi_1}]<\infty$ for some $\lambda_0>0$, gives $A_\lambda<1$ for sufficiently small $\lambda>0$. If fresh reproduction laws are independent across generations and trials, as must still be verified for the restart construction, then the correct one-trial transform is
    }{
        有 $A_0=1-p_0<1$. 由某个 $\lambda_0>0$ 满足 $\mathbb E[e^{\lambda_0\Xi_1}]<\infty$, 控制收敛定理给出足够小的 $\lambda>0$ 时 $A_\lambda<1$. 若各代和各次试验的新繁殖律独立 (对重启构造仍需验证), 则一次试验的正确变换为
    }
    \begin{equation*}
        \mathbb E[e^{\lambda X_k}]
        =\sum_{g=0}^{\infty}p_0A_\lambda^g
        =\frac{p_0}{1-A_\lambda}
        =:m_\lambda<\infty.
    \end{equation*}
    \bilingual{
        The value $m_\lambda$ need not be at least one: the conditional displacement during non-splitting generations may have negative mean. Thus no submartingale assertion is needed. Extending the trial increments $X_k$ beyond $\kappa$ with fresh laws and verifying their independence would give $\mathbb E[e^{\lambda S_k}]=m_\lambda^k$. A union bound and Markov's inequality then yield
    }{
        $m_\lambda$ 未必大于等于 $1$: 未分叉代的条件位移均值可能为负, 故不需要断言下鞅性质. 若用新繁殖律将 $X_k$ 扩展到 $k>\kappa$ 并验证独立性, 则 $\mathbb E[e^{\lambda S_k}]=m_\lambda^k$. 并集界与 Markov 不等式给出
    }
    \begin{align*}
        \mathbb P\left(
            \max\mathcal Z^{\mathrm b}_\tau\ge x_N,\,\kappa\le K
        \right)
        &\le
        \mathbb P\left(\max_{1\le k\le K}e^{\lambda S_k}
            \ge e^{\lambda c\delta\log N}\right)
        \\ &\le
        e^{-\lambda c\delta\log N}\sum_{k=1}^K m_\lambda^k
        \\ &\le
        K e^{-\lambda c\delta\log N}
            \max\{1,m_\lambda^K\}
        \le N^{-\varepsilon_3}.
    \end{align*}
    \bilingual{
        Choose $c>0$ so that $\lambda c>\max\{0,\log m_\lambda\}$. Since $K=\lfloor\delta\log N\rfloor$, the factor $K=O(\log N)$ is absorbed by any $\varepsilon_3\in(0,\delta(\lambda c-\max\{0,\log m_\lambda\}))$ for large $N$. The model-specific independence needed above remains an explicit proof obligation.
    }{
        取 $c>0$ 使 $\lambda c>\max\{0,\log m_\lambda\}$. 由于 $K=\lfloor\delta\log N\rfloor$, 对任意 $\varepsilon_3\in(0,\delta(\lambda c-\max\{0,\log m_\lambda\}))$, 因子 $K=O(\log N)$ 在 $N$ 足够大时可被幂次吸收. 以上所需的模型特定独立性仍是明确的待证条件.
    }

    \bilingual{
        We are left with the right part of \eqref{eq:a:0:2}. Since
    }{
        还需证明 \eqref{eq:a:0:2} 中的右部分. 由于
    }
    \begin{align*}
        x_N
        &=
        c\delta\log N
        +
        (\kappa-1)M
        +
        \left(
            \max\mathcal Z^{\mathrm b}_\tau
            -
            \mathcal V^{\mathrm b}(\rho_\kappa)
        \right)
        \\ &
        \le
        c\delta\log N
        +
        KM
        +
        \ell M
        \le
        \left((c+M)\delta+\frac M{\log m}\alpha\right)\log N,
    \end{align*}
    \bilingual{
        we have, since $\tau$ is a stopping time, by \nameref{lemma:branching-property-stopping-time} and the construction of $\mathcal V^{\mathrm b}$,
    }{
        由于$\tau$是停时, 根据\nameref{lemma:branching-property-stopping-time}以及 $\mathcal V^{\mathrm b}$ 的构造, 有
    }
    \begin{align*}
        &\quad
        \mathbb P\left(
            \mathcal M^{\mathrm b}_{\ell_N}
            -
            \max\mathcal Z^{\mathrm b}_\tau
            \ge
            \zeta'\log N-x_N,\,\kappa\le K,\,\tau\le T
        \right)
        \\
        &
        \le
        \mathbb P\left(
            \mathcal M^{\mathrm b}_{\ell_N}
            -
            \max\mathcal Z^{\mathrm b}_\tau
            \ge
            \left(\zeta'-(c+M)\delta-\frac M{\log m}\alpha\right)\log N,\,\kappa\le K,\,\tau\le T
        \right)
        \\ &
        =
        \mathbb E\left[
            \left.
                \mathbb P\left(
                    \mathcal M^{N|\alpha}_{\ell_N-\iota}
                    \ge
                    \left(\zeta'-(c+M)\delta-\frac M{\log m}\alpha\right)\log N
                \right)
            \right|_{\iota=\tau}
            {\bf 1}_{\{\kappa\le K,\,\tau\le T\}}
        \right]
        \\ &
        \le
        \mathbb E\left[
            N^{-\varepsilon_4}
            {\bf 1}_{\{\kappa\le K,\,\tau\le T\}}
        \right]
        \le
        N^{-\varepsilon_4},
    \end{align*}
    \bilingual{
        where $\ell_N-\iota$ is still of order $t\log^3N+o_N(\log^3N)$ when $\iota=\tau\le T=\gamma\log N$, and we can apply \eqref{eq:4.15} for $a=\alpha$ to get the second last inequality with some $\varepsilon_4=\varepsilon_4(\alpha,\delta)>0$ and some $N_4=N_4(\alpha,\delta)>0$ and all $N\ge N_4$, so long as we choose $\delta>0$ sufficiently small and $\alpha\in(0,1)$ sufficiently small such that
    }{
        其中, 当$\iota=\tau\le T=\gamma\log N$时, $\ell_N-\iota$ 仍然是 $t\log^3N+o_N(\log^3N)$ 的数量级, 因此可对 $a=\alpha$ 应用 \eqref{eq:4.15} 来得到倒数第二个不等式, 它对某个 $\varepsilon_4=\varepsilon_4(\alpha,\delta)>0$ 和某个 $N_4=N_4(\alpha,\delta)>0$ 以及所有 $N\ge N_4$ 成立, 只要我们选择 $\delta>0$ 足够小且 $\alpha\in(0,1)$ 足够小以使得
    }
    \begin{equation*}
        \zeta'-(c+M)\delta-\frac M{\log m}\alpha
        >
        -\alpha+\frac{\pi^2}2\sigma^2 t,
    \end{equation*}
    \bilingual{
        for the fixed $\zeta'>\frac{\pi^2}2\sigma^2 t$.
    }{
        对于固定的 $\zeta'>\frac{\pi^2}2\sigma^2 t$成立.
    }

    \bilingual{
        Finally, for some $\varepsilon:=\frac 12\min_{1\le i\le 4}\varepsilon_i$, some $N_0>\max_{1\le i\le 4}N_i$ and all $N\ge N_0$,
    }{
        最后, 对于某个 $\varepsilon:=\frac 12\min_{1\le i\le 4}\varepsilon_i$, 某个 $N_0>\max_{1\le i\le 4}N_i$ 以及所有 $N\ge N_0$,
    }
    \begin{align*}
        \mathbb P\left(
            \mathcal M^{N|0}_{\ell_N}
            \ge
            \zeta'\log N
        \right)
        &
        \le
        \sum_{i=1}^4N^{-\varepsilon_i}
        \le
        N^{-\varepsilon}.
    \end{align*}
    \bilingual{
        We are now done with the proof of \eqref{eq:4.15} for all $a\in[0,1]$.
    }{
        至此, 对所有 $a\in[0,1]$ 的情形, \eqref{eq:4.15} 得证.
    }
\end{proof}
